\documentclass[11pt,a4paper]{article}

\usepackage[thesenpapier]{../../templates/clifford-seminar}
\usepackage{tikz-cd}

% -----------------------------------------------------------------------------
% Document Metadata
% -----------------------------------------------------------------------------
\seminartitle[Universal Clifford Construction and Involutions]{%
  Universal Clifford Construction, Geometric Product, and Involutions%
}
\seminarsubtitle{Topic~2 --- Part~A: Foundations and Geometric Analysis}
\seminarsession{Session 02}
\seminarmodule{Part A: Foundations and Geometric Analysis}
\seminardate{November 2026}
\seminarinstitution{Universidad Distrital Francisco Jos\'e de Caldas}

\seminarauthors{Julian L. Avila-Martinez \and Cristian Martinez}
\seminarreferat{Julian L. Avila-Martinez \& Cristian Martinez}

\begin{document}

\makeseminarheader

\begin{seminarabstract}[Conceptual Scope]
  Session~01 showed that a real linear space $V$ needs a quadratic form
  $q : V \to \R$ for measurement, while its tensor algebra $T(V)$ carries
  unconstrained multilinear excess.
  In this single-session meeting we quotient $T(V)$ by
  $\mathcal{I}_q = ( v \otimes v - q(v)1 )$ to build the universal
  Clifford algebra $\Cl(V,q)$, unpack $uv = u \lc v + u \wedge v$ via
  $\Cl_{1,0}$ and $\Cl_{0,1}$, derive the $\Z_2$-grading and three canonical
  involutions from universality, and use involution parity to induce quadratic
  norms---anchored throughout in planar Kepler kinematics.
\end{seminarabstract}

% =============================================================================
\section{Motivation: Factoring Measurement into Keplerian Multiplication}
\label{sec:motivation}
% =============================================================================

In a real quadratic space $(V, q)$ with polarized symmetric bilinear form
$b(u, v) = \tfrac{1}{2}\bigl(q(u + v) - q(u) - q(v)\bigr)$, the free tensor
algebra $T(V) = \bigoplus_{r=0}^{\infty} T^r(V)$ lets us form $r \otimes p$
from the position $r$ and momentum $p$ of a Keplerian particle, yet remains
blind to $q$: $r \otimes r$ does not reduce to the squared radius $q(r) =
|r|^2$, and $r \otimes p$ encodes neither radial contraction $r \lc p$ nor the
oriented orbital plane $L = r \wedge p$
\cite{hestenesNewFoundationsClassical2002,doranGeometricAlgebraPhysicists2013}.
Internalizing $q$ into an associative product means enforcing $v^2 = q(v) 1$
for all $v \in V$---algebraically, taking the quotient of $T(V)$ by the
smallest two-sided ideal containing every measurement defect
$v \otimes v - q(v)1$ \cite{vazIntroductionCliffordAlgebras2016,
garlingCliffordAlgebrasIntroduction2011}.

% =============================================================================
\section{Foundational Definitions and Core Results}
\label{sec:definitions}
% =============================================================================

Throughout, $(V, q)$ is an $n$-dimensional real quadratic space with symmetric
bilinear form $b$, and $T(V) = \bigoplus_{r=0}^{\infty} T^r(V)$ is its unital
tensor algebra ($1 \in T^0(V) = \R$).

\subsection{Ideals, Quotient Algebra, and the Universal Property}
\label{subsec:quotient-universal}

\begin{definition}[Two-Sided Ideal and Universal Clifford Algebra]
  \label{def:clifford-algebra}
  A subspace $\mathcal{I} \subseteq A$ of an associative unital $\R$-algebra
  $A$ is a \emph{two-sided ideal} if $ax, xa \in \mathcal{I}$ for all
  $a \in A$, $x \in \mathcal{I}$; the quotient space $A/\mathcal{I}$
  \cite[Sec.~3.E]{axlerLinearAlgebraDone2015} then
  inherits a well-defined associative unital multiplication $[a][b] := [ab]$
  and canonical surjection $\pi : A \to A/\mathcal{I}$ with $\ker\pi =
  \mathcal{I}$.
  Setting
  \begin{equation}
    \label{eq:clifford-ideal}
    \mathcal{I}_q := \left( v \otimes v - q(v) 1 \;\middle|\; v \in V
    \right) \subset T(V), \qquad
    \Cl(V, q) := T(V) / \mathcal{I}_q,
  \end{equation}
  defines the \emph{universal Clifford algebra} of $(V, q)$, equipped with the
  canonical Clifford map $\iota := \pi \circ j : V \hookrightarrow T(V) \to
  \Cl(V, q)$.
  Identifying $V$ with its injective image $\iota(V)$ and writing the quotient
  product by juxtaposition yields $v^2 = q(v) 1$ for all $v \in V$.
\end{definition}

A linear map $\gamma : V \to A$ into an associative unital $\R$-algebra $A$ is
a \emph{Clifford map} if $\gamma(v)^2 = q(v)1_A$ for all $v \in V$
(equivalently, $\gamma(u)\gamma(v) + \gamma(v)\gamma(u) = 2b(u, v)1_A$).

\begin{theorem}[Universal Property of $\Cl(V, q)$]
  \label{thm:universal-property}
  For every associative unital $\R$-algebra $A$ and every Clifford map
  $\gamma : V \to A$, there exists a unique unital algebra homomorphism
  $\Phi : \Cl(V, q) \to A$ such that $\Phi \circ \iota = \gamma$:
  \begin{equation}
    \label{eq:universal-diagram}
    \begin{tikzcd}[column sep=large, row sep=small]
      V \arrow[r, "\iota"] \arrow[dr, "\gamma"'] & \Cl(V, q) \arrow[d,
        dashed, "\exists!\,\Phi"] \\
      & A
    \end{tikzcd}
  \end{equation}
  Hence $\Cl(V, q)$ is unique up to a unique isomorphism and has dimension
  $\dim \Cl(V, q) = 2^n$
  \cite[Secs.~3.1--3.2]{vazIntroductionCliffordAlgebras2016}.
\end{theorem}

\subsection{The Geometric Product and One-Dimensional Signatures}
\label{subsec:geom-product}

\begin{proposition}[Polarization and Vector Invertibility]
  \label{prop:polarization}
  Any $u, v \in V \subset \Cl(V, q)$ satisfy $uv + vu = 2b(u, v) 1$.
  In particular, $u \perp_b v \iff uv = -vu$, and every non-isotropic vector
  $v \in V$ ($q(v) \neq 0$) admits the unique two-sided inverse
  $v^{-1} = v / q(v)$.
\end{proposition}

\begin{keyresult}[Decomposition of the Geometric Product]
  For all $u, v \in V \subset \Cl(V, q)$, the associative geometric product
  splits uniquely into its symmetric contraction and antisymmetric exterior
  parts:
  \begin{equation}
    \label{eq:geometric-product}
    uv = u \lc v + u \wedge v, \qquad
    u \lc v = u \rc v := \frac{1}{2}(uv + vu) = b(u, v) 1, \quad
    u \wedge v := \frac{1}{2}(uv - vu),
  \end{equation}
  where $\lc$ and $\rc$ denote the \emph{left} and \emph{right contractions}
  (coinciding on $1$-vectors and generalizing to $A_r \lc B_s =
  \grade{A_r B_s}{s-r}$ and $A_r \rc B_s = \grade{A_r B_s}{r-s}$).
  In the Kepler plane, $rp = r \lc p + L$ unifies radial contraction
  ($r \lc p$) and the orbital angular momentum bivector ($L = r \wedge p$),
  while vector inversion yields $r^{-1} = r / r^2$ in
  $\dot{p} = -k\, r^{-1}/|r|$.
\end{keyresult}

Already in dimension $n = 1$, signature flips the algebraic type completely:
$\Cl_{1,0} = \Cl(\R^{1,0})$ ($e_1^2 = +1$) versus $\Cl_{0,1} = \Cl(\R^{0,1})$
($e_1^2 = -1$), compared in \cref{tab:low-dim}.

\begin{table}[htbp]
  \centering
  \caption{One-dimensional real Clifford algebras $\Cl_{1,0}$ and $\Cl_{0,1}$
    for $u = a + b e_1$ ($a, b \in \R$).}
  \label{tab:low-dim}
  \small
  \renewcommand{\arraystretch}{1.16}
  \begin{tabularx}{\textwidth}{@{} >{\bfseries}l l l >{\raggedright\arraybackslash}X @{}}
    \toprule
    \textbf{Algebra} & \textbf{Square} & \textbf{Isomorphism} &
      \textbf{Structure \& Geometric Role} \\
    \midrule
    $\Cl_{1,0}$ & $e_1^2 = +1$ &
      $\R[x]/(x^2 - 1) \cong \R \oplus \R$ &
      Split-complex (hyperbolic) numbers; non-simple with idempotents
      $P_\pm = \tfrac{1}{2}(1 \pm e_1)$ ($P_+ P_- = 0$, zero divisors) and
      $u\clconj{u} = a^2 - b^2$. \\
    $\Cl_{0,1}$ & $e_1^2 = -1$ &
      $\R[x]/(x^2 + 1) \cong \C$ &
      Complex field (simple division algebra); $e_1$ acts as $i$ and models the
      even rotation subalgebra $\Cl^+_{2,0} \cong \C$ of planar Kepler orbits
      ($e^{\theta \hat{B}}$). \\
    \bottomrule
  \end{tabularx}
\end{table}

\subsection{\texorpdfstring{$\Z_2$}{Z2}-Grading, Involutions, and Parity-Induced Norms}
\label{subsec:grading-involutions-norm}

Because $v^2 = q(v) 1$ mixes tensor degrees $2$ and $0$, $\Cl(V, q)$ retains
only a $\Z_2$-grading as an algebra.
All three canonical involutions follow coordinate-freely from
\cref{thm:universal-property}
\cite[Sec.~5.3]{garlingCliffordAlgebrasIntroduction2011}:
(i)~\textbf{grade involution} $\ginv{\cdot}$, the unique automorphism extending
$v \mapsto -v$;
(ii)~\textbf{reversion} $\rev{\cdot}$, the unique anti-automorphism
($\rev{AB} = \rev{B}\rev{A}$) induced by $v \mapsto v$ into the opposite
algebra $\Cl(V,q)^{\mathrm{opp}}$; and
(iii)~\textbf{Clifford conjugation} $\clconj{A} := \ginv{\rev{A}} =
\rev{\ginv{A}}$.
The $\pm 1$ eigenspaces $\Cl^\pm(V, q) := \{ A \mid \ginv{A} = \pm A \}$ yield
the superalgebra decomposition $\Cl(V, q) = \Cl^+ \oplus \Cl^-$ with
\begin{equation}
  \label{eq:z2-grading}
  \Cl^\pm \Cl^\pm \subseteq \Cl^+, \qquad
  \Cl^\pm \Cl^\mp \subseteq \Cl^-.
\end{equation}
Thus $\Cl^+(V, q)$ is a closed subalgebra (containing scalars, bivectors such
as $L$, and rotors), while vectors ($r, p$) lie in $\Cl^-(V, q)$.
Their action on a homogeneous $k$-blade $A_k \in \Cl^k(V, q) \cong
\bigwedge^k V$ is periodic modulo~$4$ (\cref{tab:involutions}).

\begin{table}[htbp]
  \centering
  \caption{Canonical involutions on a $k$-blade $A_k \in \Cl^k(V, q)$
    (periodic modulo~$4$).}
  \label{tab:involutions}
  \small
  \renewcommand{\arraystretch}{1.15}
  \begin{tabularx}{\textwidth}{@{} >{\bfseries}l c c c c c c >{\raggedright\arraybackslash}X @{}}
    \toprule
    \textbf{Involution} & \textbf{Map} & \textbf{Factor} &
      $k{=}0$ & $k{=}1$ & $k{=}2$ & $k{=}3$ & \textbf{Algebraic Character} \\
    \midrule
    Grade involution & $\ginv{A}_k$ & $(-1)^k$ & $+$ & $-$ & $+$ & $-$ &
      Automorphism ($\Z_2$-parity) \\
    Reversion & $\rev{A}_k$ & $(-1)^{k(k-1)/2}$ & $+$ & $+$ & $-$ & $-$ &
      Anti-automorphism (order swap) \\
    Clifford conjugation & $\clconj{A}_k$ & $(-1)^{k(k+1)/2}$ & $+$ & $-$ & $-$ & $+$ &
      Anti-automorphism ($\clconj{A} = \ginv{\rev{A}}$) \\
    \bottomrule
  \end{tabularx}
\end{table}

\begin{keyresult}[Involution Parity and Induced Quadratic Norms]
  Let $\grade{\cdot} : \Cl(V, q) \to \R$ denote the scalar projection (grade-$0$
  part, omitting the subscript when $k = 0$).
  The two anti-automorphisms induce canonical quadratic forms on $\Cl(V, q)$
  \cite[Sec.~3.3]{vazIntroductionCliffordAlgebras2016}:
  \begin{equation}
    \label{eq:involution-norms}
    N_{\sim}(A) := \grade{A\rev{A}}, \qquad
    N_{-}(A) := \grade{A\clconj{A}}, \qquad
    N_{-}(A_k) = (-1)^k N_{\sim}(A_k),
  \end{equation}
  linked on every $k$-blade by the $\Z_2$-parity sign $(-1)^k$ from
  $\clconj{A}_k = \ginv{\rev{A}}_k$.
  Hence $N_\sim$ is positive-definite on Euclidean $\Cl_{n,0}$
  ($e_i\rev{e}_i = +1$), while $N_-$ is positive-definite on $\Cl_{0,n}$
  ($e_i\clconj{e}_i = -e_i^2 = +1$, recovering $|z|^2 = z\clconj{z}$ in
  $\Cl_{0,1} \cong \C$).
\end{keyresult}

\begin{remark}[Keplerian Magnitudes]
  In $\Cl_{2,0} \subset \Cl_{3,0}$, $r$ is odd ($k=1$), so $r\rev{r} = |r|^2 >
  0$ and $r^{-1} = \rev{r}/(r\rev{r})$.
  By contrast, $L = r \wedge p$ is an even bivector ($k=2$) with $L^2 =
  -|L|^2 \le 0$; reversion supplies the minus sign yielding
  $L\rev{L} = L\clconj{L} = -L^2 = |L|^2 = (2m\dot{\mathcal{A}})^2$, the
  squared Keplerian areal velocity.
\end{remark}

% =============================================================================
\section{Blackboard Derivation Agenda (40 Minutes)}
\label{sec:agenda}
% =============================================================================

\begin{enumerate}
  \item \textbf{Stage 1 (12 min):} Ideal $\mathcal{I}_q$, quotient
    $\Cl(V, q) = T(V)/\mathcal{I}_q$, and proof of the universal property
    (\cref{thm:universal-property}).
  \item \textbf{Stage 2 (10 min):} Geometric product $uv = u \lc v +
    u \wedge v$, contractions $\lc,\rc$, vector inverses, and $\Cl_{1,0}$,
    $\Cl_{0,1}$.
  \item \textbf{Stage 3 (9 min):} Coordinate-free construction of
    $\ginv{\cdot}$, $\rev{\cdot}$, $\clconj{\cdot}$, $\Z_2$-grading, and
    $k$-blade signs.
  \item \textbf{Stage 4 (9 min):} Parity-induced norms $N_\sim$ vs.\ $N_-$ on
    $\Cl_{1,0}$, $\Cl_{0,1}$, and the Kepler product $rp = r \lc p + L$.
\end{enumerate}

% =============================================================================
\section{Theses for Plenary Discussion}
\label{sec:theses}
% =============================================================================

\begin{thesis}[Metric Quotients Over Free Multilinearity]
  Factoring $T(V)$ by $\mathcal{I}_q = ( v \otimes v - q(v)1 )$ is
  the minimal algebraic act turning a linear space into a metric algebra; by
  universality, $\C$, $\mathbb{H}$, and Pauli matrices are homomorphic shadows
  of a single invariant structure.
\end{thesis}

\begin{thesis}[Signature Governs Divisibility and Geometry]
  Even in dimension one, reversing the sign of $q$ transforms a division
  algebra of rotations ($\Cl_{0,1} \cong \C$) into a non-simple ring with
  zero divisors and idempotents ($\Cl_{1,0} \cong \R \oplus \R$); zero divisors
  are not defects, but the algebraic signature of hyperbolic and null geometry.
\end{thesis}

\begin{thesis}[Involution Parity Selects Physical Norms]
  Reversion and Clifford conjugation differ strictly by the $\Z_2$-parity
  automorphism $\ginv{\cdot}$, selecting the positive-definite norm for each
  signature ($|z|^2 = z\clconj{z}$ in $\Cl_{0,1}$; $|r|^2 = r\rev{r}$ and
  $|L|^2 = L\rev{L}$ in $\Cl_{3,0}$) without coordinates.
\end{thesis}

% =============================================================================
\section{Preparatory Exercises}
\label{sec:exercises}
% =============================================================================

Attempt these three exercises and submit one written question 24~hours before
the session.

\begin{enumerate}[itemsep=0.2em, topsep=0.3em]
  \item \textbf{One-Dimensional Quotients and Zero Divisors.}
    Identify $T(\R) \cong \R[x]$ and construct $\Cl_{1,0} \cong
    \R[x]/(x^2 - 1)$ and $\Cl_{0,1} \cong
    \R[x]/(x^2 + 1)$.
    In $\Cl_{1,0}$, set $P_\pm := \tfrac{1}{2}(1 \pm e_1)$ and verify
    $P_\pm^2 = P_\pm$, $P_+ + P_- = 1$, and $P_+ P_- = 0$, proving
    $\Cl_{1,0} \cong \R \oplus \R$.
    Explain why $\R^{1,0}$ admits a one-dimensional non-universal Clifford
    algebra whereas $\R^{0,1}$ does not.

  \item \textbf{Coordinate-Free Grade Involution.}
    Using only \cref{thm:universal-property} (no basis), prove that
    $v \mapsto -\iota(v)$ extends uniquely to an involutory algebra
    automorphism $\ginv{\cdot} : \Cl(V, q) \to \Cl(V, q)$, and deduce directly
    that $\Cl^+(V, q) = \{ A \mid \ginv{A} = A \}$ is a subalgebra.

  \item \textbf{Geometric Product and Parity-Induced Norms in Kepler Motion.}
    \begin{enumerate}
      \item For $r, p \in \R^{2,0} \subset \Cl_{2,0}$ with
        $rp = r \lc p + L$ ($L = r \wedge p$), verify $\rev{L} = -L$ and
        compute $(rp)\rev{(rp)}$ to deduce the Lagrange identity
        $(r \lc p)^2 + |L|^2 = r^2 p^2$, where $|L|^2 := L\rev{L} = -L^2$.
      \item For $u = a + b e_1$ ($a, b \in \R$), compute $u\rev{u}$,
        $u\clconj{u}$, $N_\sim(u)$, and $N_-(u)$ in $\Cl_{1,0}$ and
        $\Cl_{0,1}$, and explain why $\ginv{e}_1 = -e_1$ makes $u\clconj{u}$ a
        scalar multiplicative norm while $u\rev{u}$ is not when $ab \neq 0$.
    \end{enumerate}
\end{enumerate}

% =============================================================================
\section{Preparatory Literature}
\label{sec:literature}
% =============================================================================

\begin{itemize}[itemsep=0.15em, topsep=0.25em]
  \item \textbf{Primary reading:} Vaz and da Rocha
    \cite[Secs.~3.1--3.3]{vazIntroductionCliffordAlgebras2016} (quotient,
    universality, $\Z_2$-grading, involutions, norms); Axler
    \cite[Sec.~3.E]{axlerLinearAlgebraDone2015} (quotient linear spaces).
  \item \textbf{Further references:} Garling
    \cite[Secs.~5.1--5.3, 5.6]{garlingCliffordAlgebrasIntroduction2011};
    Lounesto \cite[Chaps.~1--3]{lounestoCliffordAlgebrasSpinors2001};
    Hestenes \cite[Chaps.~1--2, Sec.~4-1]{hestenesNewFoundationsClassical2002};
    Doran and Lasenby \cite[Chaps.~1--2]{doranGeometricAlgebraPhysicists2013}.
\end{itemize}

{\renewcommand*{\bibfont}{\footnotesize}\setlength{\bibitemsep}{0.15\baselineskip}%
\printbibliography[title={References}]}

\end{document}
